\documentclass{article}
\usepackage[T1]{fontenc}
\usepackage[english,polish]{babel}
\usepackage{geometry}
\usepackage{datetime}
\geometry{a4paper, margin=1in}
\usepackage{fancyhdr}
\usepackage{tabularray}
\usepackage{hyperref}

\hypersetup{
    colorlinks=true,
    linkcolor=blue,
    filecolor=green,
    urlcolor=blue
}
\pagestyle{fancy}
\setlength{\headheight}{35pt}
\fancyhf{}
\fancyhead[L]{Systemy wbudowane i mikroprocesory}
\fancyhead[R]{Jakub Chyła\\Projekt\\25.04.2026}
\renewcommand{\today}{\number\year-\number\month-\number\day}

\AtBeginDocument{\selectlanguage{polish}}

\begin{document}

\section{Opis}
\begin{flushleft}
Celem projektu jest stworzenie systemu, który zapewni utrzymanie odpowiedniej wilgotności gleby dla roślin doniczkowych.
W tym celu wykorzystano pompę do pobierania wody oraz czujnik wilgotności gleby.\\
\end{flushleft}
\section{Sprzęt}
\begin{itemize}
    \item Płytka deweloperska STM32 L412KB
    \item Wyświetlacz SH1106 128x64 (I2C)
    \item Pojemnościowy czujnik wilgotności v1.2
    \item Przekaźnik 4-kanałowy 12V z wbudowanym optoizolatorem
    \item Pompa wodna 12V 0.26A
    \item Zasilacz 12V/1A
    \item Zasilacz Micro-USB
\end{itemize}
\begin{flushleft}
    Do zasilenia płytki, czujnika wilgotności, ekranu oraz części logicznej układu przekaźnika wykorzystano zasilacz Micro-USB, pompa oraz część mechaniczna przekaźnika
    zasilane są zasilaczem 12V. Z powodu zbyt małego napięcia wyjściowego na pinach płytki STM32L412KB (3,3V) i żeby zabezpieczyć płytkę, wykorzystano optoizolator wbudowany
    w układ przekaźnika.\\
    Przeprowadzono kalibrację czujnika poprzez sprawdzenie podawanych danych w suchym powietrzu oraz po zanurzeniu w wodzie. Wartość dla powietrza wynosiła około 2700,
    natomiast po zanurzeniu w wodzie około 1450.
\end{flushleft}
\begin{center}
    \begin{tblr}{colspec={cc}, hlines, vlines}
        STM32 Pin & Urządzenie \\
        D0 & SH1106 SDA\\
        D1 & SH1106 SCK\\
        D5 & Przekaźnik IN1\\
        A0 & Sensor AOUT\\
        3V3 & {SH1106 VDD\\Przekaźnik VCC\\Sensor VCC}\\
        GND & {SH1106 GND\\Sensor GND}\\
    \end{tblr}
\end{center}
\section{Kod}
\href{https://github.com/jakubchyla/irrigation}{https://github.com/jakubchyla/irrigation}
\subsection{Wyświetlacz SH1106}
\begin{flushleft}
    W celu obsłużenia wyświetlacza SH1106 napisano bibliotekę, która umożliwia inicjalizację oraz wyświetlanie tekstu pod danym adresem na ekranie.
    Do komunikacji z wyświetlaczem wykorzystywany jest protokół I2C.\\
    \vspace{7pt}
    \href{https://github.com/jakubchyla/irrigation/blob/master/Core/Src/sh1106.c}{Github - sh1106.c}\\
    \href{https://github.com/jakubchyla/irrigation/blob/master/Core/Inc/sh1106.h}{Github - sh1106.h}

\end{flushleft}
\subsection{Czujnik wilgotności}
\begin{flushleft}
    Czujnik zwraca dane w postaci analogowej, do konwersji wykorzystano wbudowany w płytkę ADC, a następnie przetworzono otrzymany wynik
    na skalę 0-100 na podstawie danych z kalibracji. Wilgotność poniżej, której uruchamiana jest pompa można skonfigurować poprzez
    zmianę makr SENSOR<numer>\_TARGET\\
    \vspace{7pt}
    \href{https://github.com/jakubchyla/irrigation/blob/master/Core/Src/irrigation.c}{Github - irrigation.c}\\
    \href{https://github.com/jakubchyla/irrigation/blob/master/Core/Inc/irrigation.h}{Github - irrigation.h}
\end{flushleft}

\subsection{Przekaźnik, pompa}
\begin{flushleft}
    Pompa jest uruchamiana poprzez wysłanie sygnału LOW na dany pin (IN1/IN2/IN3/IN4) przekaźnika.
    Czas pracy pompy można skonfigurować poprzez zmianę makr PUMP<numer>\_TIME.\\
    \vspace{7pt}
    \href{https://github.com/jakubchyla/irrigation/blob/master/Core/Src/irrigation.c}{Github - irrigation.c}\\
    \href{https://github.com/jakubchyla/irrigation/blob/master/Core/Inc/irrigation.h}{Github - irrigation.h}
\end{flushleft}

\end{document}
